\documentclass[12pt]{article}
\usepackage{amsmath, amssymb, amsfonts}
\usepackage[T1]{fontenc}
\usepackage[utf8]{inputenc}
\usepackage{geometry}
\geometry{margin=2.5cm}
\usepackage{lmodern}

\title{Ecuaciones del Indicador IS$_d$\\
\large Indicador de Singularidad Documental}
\author{Blázquez Ochando, M. · Ovalle Perandones, M.A. · Prieto Gutiérrez, J.J.\\
Universidad Complutense de Madrid}
\date{2025}

\begin{document}
\maketitle

% ────────────────────────────────────────────────────────────────
\section*{Ecuación (1) — Estructura general del indicador}
% ────────────────────────────────────────────────────────────────

\begin{equation}
IS_d \;=\; \alpha \cdot O_d \;+\; (1-\alpha) \cdot I_d
\qquad \alpha \in [0,1]
\label{eq:isd}
\end{equation}

\noindent
Donde $O_d$ es el componente de originalidad textual, $I_d$ el componente de
impacto bibliométrico y $\alpha$ el parámetro de ponderación.

% ────────────────────────────────────────────────────────────────
\section*{Ecuación (2) — Normalización de $O_d$}
% ────────────────────────────────────────────────────────────────

\begin{equation}
O_d \;=\; \min\!\left(\frac{O_d^{raw}}{P_{99}(O^{raw})},\; 1\right)
\label{eq:od}
\end{equation}

\noindent
Donde $P_{99}(O^{raw})$ es el percentil 99 de los valores $O_d^{raw}$ en el corpus.

% ────────────────────────────────────────────────────────────────
\section*{Ecuación (3) — Originalidad textual bruta $O_d^{raw}$}
% ────────────────────────────────────────────────────────────────

\begin{equation}
O_d^{raw} \;=\;
\frac{\displaystyle\sum_{i \,\in\, \Omega_d} S_i \cdot \log\!\left(1 + f_i^{\,doc}\right)}
     {\;\left|\Omega_d\right| \cdot \log\!\!\left(1 + \dfrac{T_d}{\left|\Omega_d\right|}\right) + \varepsilon\;}
\label{eq:odraw}
\end{equation}

\noindent
Donde:
\begin{itemize}
  \item $\Omega_d$ = conjunto de n-gramas únicos presentes en el documento $d$
  \item $|\Omega_d|$ = número de n-gramas únicos (cardinal del conjunto)
  \item $f_i^{\,doc}$ = frecuencia absoluta del n-grama $i$ en el documento $d$
  \item $T_d$ = total de n-gramas extraídos del documento $d$ (con repetición)
  \item $\varepsilon = 10^{-6}$ = constante de suavizado numérico
\end{itemize}

% ────────────────────────────────────────────────────────────────
\section*{Ecuación (4) — Peso de singularidad $S_i$}
% ────────────────────────────────────────────────────────────────

\begin{equation}
S_i \;=\;
\underbrace{\log\!\left(\frac{C+1}{cd_i + 0{,}5}\right)}_{\text{IDF de Robertson}}
\;\cdot\;
\underbrace{\frac{1}{1 + \log\!\left(1 + f_i^{\,corpus}\right)}}_{\text{penalización de frecuencia (BM25)}}
\label{eq:si}
\end{equation}

\noindent
Donde:
\begin{itemize}
  \item $C$ = número total de documentos en el corpus
  \item $cd_i$ = número de documentos del corpus que contienen el n-grama $i$
  \item $f_i^{\,corpus}$ = frecuencia total del n-grama $i$ en todo el corpus
\end{itemize}

% ────────────────────────────────────────────────────────────────
\section*{Ecuación (5) — Impacto bibliométrico normalizado $I_d$}
% ────────────────────────────────────────────────────────────────

\begin{equation}
I_d \;=\; \frac{\log(1 + C_d)}{\log(1 + C_{max})}
\label{eq:id}
\end{equation}

\noindent
Donde $C_d$ es el número de citas recibidas por el documento $d$ y
$C_{max}$ es el máximo de citas en el corpus ($C_{max} = 53{.}982$ en calcia.db).
Se verifica que $I_d \in [0,1]$ exactamente: $I_d = 0$ si $C_d = 0$; $I_d = 1$
si $C_d = C_{max}$.

% ────────────────────────────────────────────────────────────────
\section*{Ecuación ($\alpha^*$) — Criterio objetivo de calibración}
% ────────────────────────────────────────────────────────────────

\begin{equation}
\alpha^* \;=\; \arg\min_{\alpha \,\in\, [0,1]}
\left|\,\mathrm{skew}\!\left(\alpha \cdot O_d + (1-\alpha)\cdot I_d\right)\right|
\label{eq:alpha}
\end{equation}

\noindent
Sobre calcia.db: $\alpha^* = 0{,}692 \approx 0{,}70$, con $|\mathrm{skew}| = 0{,}00058$
(mínimo global). La T de Kendall de los deciles permanece igual a 1,0 para
todo $\alpha \in [0;\, 0{,}90]$.

\end{document}